\chapter*{Abstract}

In this thesis, we developed a Finite Element Method (FEM)-anchored
Physics-Informed Neural Network (PINN) framework for transient water-quenching
simulation of A356 aluminium.
We based the study on the industrial quenching scenario of Mortensen
et al.~\cite{mortensen2026}.
Our central question was: how many FEM solver calls can be replaced by PINN
predictions while keeping the temperature error within engineering limits?

We adopted Multi-Step Window Prediction (MSWP) as the core strategy.
The FEM solver provides anchor fields at selected time steps.
The PINN predicts between consecutive anchors.
The skip factor $k$ controls the anchor spacing.
We compared three NAS (Neural Architecture Search)-selected architectures:
Bayesian/TPE ($5{\times}151$, ReLU), NSGA-II ($3{\times}153$, tanh), and
NSGA-III ($3{\times}75$, tanh).
We conducted experiments across ten unique benchmark geometries in two studies.
The canonical study covers eight domains: three 2D, four 3D, and the Thermal
Fin, originally a 2D steady-state RBniCS benchmark that we extended to a 3D
transient quenching scenario.
A complementary adaptive-$k$ study added Square and Flower as two new 2D
geometries alongside Circle and L-Shape.

On the adaptive-$k$ study, we found 83--85\% FEM saving with mean absolute
error (MAE) below $2.5\,^\circ\mathrm{C}$ across all architectures.
On the canonical 2D domains with fixed $k$, Bayesian/TPE stayed below
$5\,^\circ\mathrm{C}$ through $k=3$ on all three domains.
On the 3D domains, Bayesian/TPE remained acceptable at $k=5$ with MAE below
$14\,^\circ\mathrm{C}$.
On the Thermal Fin, Bayesian/TPE with Fourier embedding reached $k=5$ at
$12.55\,^\circ\mathrm{C}$ (ten-seed mean).
NSGA-II and NSGA-III held through $k=4$.
We demonstrated that FEM-anchored PINNs can reproduce transient quenching
fields with 65--85\% fewer solver calls within practical engineering accuracy.

\vspace{0.15cm}
\noindent\textbf{Keywords}\quad Physics-Informed Neural Networks, Neural Architecture Search, Finite Element Method, Transient Heat Transfer, Water Quenching, Multi-Step Window Prediction, Adaptive Skip-Factor Scheduling, FEM-Anchored Surrogate Modelling, Fourier Feature Embeddings, Thermal Fin Benchmark, A356 Aluminium.
